\section{Introduction}

\subsection{Problem Background}

With the continuous advancement of global industrialization and population growth, energy consumption has been steadily increasing. Although the proportion of clean energy from solar, wind, and nuclear power has risen, fossil fuels remain the dominant energy source, leading to the exacerbation of the urban heat island effect and greenhouse effect.

Passive Daytime Radiative Cooling (PDRC) has emerged as a transformative zero-energy cooling technology. Based on the Stefan-Boltzmann Law and the 8--13 $\mu$m atmospheric transparency window, PDRC enables objects to spontaneously emit thermal radiation into outer space (approximately 3 K), achieving sub-ambient cooling without any energy input. Among candidate materials, polydimethylsiloxane (PDMS) thin films exhibit ultra-high visible transmittance and high infrared emissivity arising from Si--O--Si and Si--C bond vibrations, making them highly suitable for radiative cooling applications in energy-efficient buildings, solar cells, and personal thermal management.

However, the radiative cooling performance depends critically on film thickness and structural design, governed by thin-film optical interference effects. This paper establishes a comprehensive modeling framework based on the Transfer Matrix Method (TMM) and Kirchhoff's law, extending from single-layer emissivity analysis to multilayer structure optimization with economic feasibility assessment.

\subsection{Problem Restatement}

This paper addresses four interconnected sub-problems:

\textbf{Problem 1:} Establish a mathematical model to study the emissivity of PDMS thin films as a function of wavelength for different film thicknesses.

\textbf{Problem 2:} Develop an evaluation model to assess the radiative cooling performance of PDMS thin films with varying thicknesses, providing recommendations for technology advancement.

\textbf{Problem 3:} Select suitable materials to form multilayer film structures with PDMS, incorporating optimization design algorithms for optimal material selection and layer thicknesses.

\textbf{Problem 4:} Perform comprehensive optimal design of radiative cooling materials and structures to achieve the best possible cooling performance, followed by feasibility and cost evaluation.

\subsection{Overall Approach}

Our approach proceeds through four stages: (1) single-layer PDMS emissivity modeling via TMM with Airy summation; (2) radiative cooling performance evaluation incorporating blackbody radiation, atmospheric transmittance, and solar spectrum; (3) multilayer structure design with parameter scanning across candidate materials; and (4) comprehensive optimization with economic and feasibility assessment.
